\chapter{Modelli e UML}
\begin{boxDef}
\textbf{\textcolor{brown}{Sistema:}} è un \textbf{insieme} di componenti che interagiscono e si influenzano tra loro, percepiti come unica unità per lo scopo comune. \\
\textbf{\textcolor{brown}{Astrazione:}} \textbf{semplifica} la rappresentazione della \textbf{realtà} sia nelle generalizzazioni, sia nei processi, eliminando dettagli e identificando punti in comune.
\end{boxDef}

\section{UML}
\begin{boxDef}
\textbf{\textcolor{red}{UML}} è una famiglia di notazioni grafiche per descrivere e progettare sistemi software, soprattutto quelli \textbf{OOP}. I linguaggi di programmazione sono troppo basso livello per discutere di progettazione software, i linguaggi grafici aiutano la comprensione di gruppo.
\begin{itemize}
    \item \textbf{Sketch:} maggiormente utile per forward e reverse engineering enfatizzando la comunicazione selettiva piuttosto che completa.
    \item \textbf{Blueprint:} enfasi sulla completezza e sui dettagli, che richiedono strumenti più sofisticati di quelli per costruire gli sketch.
    \item \textbf{Linguaggio di programmazione:} disegnare dei diagrammi UML che vengono compilati in codice eseguibile.
\end{itemize}
\end{boxDef}
\begin{itemize}
    \item \textbf{Regole prescrittive:} per definire il corpo e specificare cosa è legale o illegale e quale significato è associato ad ogni frase del linguaggio.
    \item \textbf{Regole descrittive:} per decidere come il linguaggio viene effettivamente utilizzato ed è quello più utilizzato nell'ambito \textbf{UML}.
\end{itemize}
È importante ricordare che se una specifica non è inclusa nell'\textbf{UML} non è detto che non sia presente nel progetto, perché contiene solo le cose utili. Esitono diversi tipi di modelli \textbf{UML}:
\begin{itemize}
    \item \textbf{\textcolor{brown}{Class Diagram.}}
    \item \textbf{\textcolor{brown}{Sequence Diagram.}}
    \item \textbf{\textcolor{brown}{Activity Diagram.}}
    \item \textbf{\textcolor{brown}{State Diagram.}}
    \item \textbf{\textcolor{brown}{Use-Case Diagram.}}
\end{itemize}

